\chapter{Testing}
\label{chap:testing}

The validation of the emulated STM32F405 I$^2$C controller and LIS3DH accelerometer required a comprehensive testing strategy that verifies functional correctness, protocol compliance, and firmware portability. Unlike unit tests that isolate individual functions, the testing approach adopted in this thesis employs full-system integration testing: unmodified STM32CubeF4 HAL-based firmware executing within QEMU, interacting with the emulated peripherals via the emulated I$^2$C bus, and producing observable outputs via UART console redirection.

\textbf{Testing objectives:}

\begin{enumerate}
    \item \textbf{Protocol Fidelity:} Verify that I$^2$C subaddressing, auto-increment, 
    and register access sequences match the LIS3DH datasheet specifications.
    
    \item \textbf{Configuration Correctness:} Confirm that writing to control registers \\
    (CTRL\_REG1 to CTRL\_REG6) produces correct internal state transitions 
    (operating mode, full-scale range, output data rate).
    
    \item \textbf{Data Acquisition:} Validate that acceleration values injected via 
    QEMU's QOM properties propagate correctly through register reads and produce 
    accurate scaled output in milligravity units.
    
    \item \textbf{Interrupt Behavior:} Test Data-Ready (DRDY) interrupt generation 
    on INT1, GPIO EXTI handling, and latched/non-latched interrupt modes.
    
    \item \textbf{HAL Driver Compatibility:} Demonstrate that STMicroelectronics' 
    reference HAL drivers (\texttt{HAL\_I2C\_Mem\_Read()}, 
    \texttt{HAL\_I2C\_Mem\_Write()}) execute without modification.
\end{enumerate}


% ==== Test Execution Enviroment ==== %
\section{STM32CubeIDE Configuration}
\label{chap:test_enviroment}

As mentioned in the section~\ref{chap:Back_ARM_STM32CubeIDE}, all firmware development for 
this project was performed using \texttt{STM32CubeIDE} and its associated tools. This 
section describes the pre-configuration carried out in \texttt{STM32CubeIDE} prior to 
defining the functional tests.

\texttt{STM32CubeIDE} allows the creation of projects specifically targeted to 
particular microcontrollers or development boards. Using this feature, a new project 
is created using a preconfigured STM32F4-based evaluation board as the base. This board 
uses the STM32F407 microcontroller, which, as previously noted, is compatible with 
STM32F405 configurations.

Once the project is created, \texttt{STM32CubeMX} opens automatically, allowing the 
microcontroller to be configured visually as follows:

\begin{enumerate}
  \item Enable I2C1 and configure it for 100~kHz operation with 7-bit addressing. 
  Then assign the I2C1\_SCL and I2C1\_SDA lines to pins PB6 and PB7, respectively.
  
  \item Configure a UART peripheral to enable debugging.
  
  \item Reconfigure pin PA1, which is wired to the LIS3DH INT1 pin, to 
  \textbf{GPIO\_EXTI1} mode. In the configuration for this pin, select 
  ``External interrupt mode with rising edge trigger detection'' and ``No pull-up and 
  no pull-down''. Then, in the \textbf{NVIC} tab, enable the \textbf{EXTI line 1 
  interrupt} and assign it a priority of 5.
\end{enumerate}

Once these configurations are applied and saved, \texttt{CubeMX} generates the code 
required to initialize the I$^2$C, GPIO, and EXTI lines. Within this generated code, and 
before integrating the LIS3DH library or defining the tests, the following two 
modifications must be made:

\begin{enumerate}
  \setcounter{enumi}{3}
  \item In the \texttt{main()} function, locate and comment out the 
  \texttt{SystemClock\_Config()} call.
  
  \item Define the \texttt{HAL\_GPIO\_EXTI\_Callback()} function, which is executed 
  automatically when the DRDY interrupt is triggered on INT1.
\end{enumerate}

The \texttt{SystemClock\_Config()} function is commented out because the STM32F405 
model implemented in QEMU does not accurately emulate the entire clock subsystem. As 
a result, some function calls within \texttt{SystemClock\_Config()} return 
\texttt{HAL\_ERROR}, which in turn causes \texttt{Error\_Handler()} to be invoked and 
enter its infinite loop.


% ==== Test Cases & Results ==== %
\section{Test Cases}
\label{chap:test_test_case}

This section presents a series of tests and their results which, in the developer's opinion, cover all aspects necessary to validate the operation of the project, both of the I2C controller and the LIS3DH.

\subsection{Test Cases}

\subsubsection{Test Case 1: Device Initialization and WHO\_AM\_I Verification}

\textbf{Objectives:} Verify that the LIS3DH device responds to I$^2$C transactions and 
returns the correct device ID.

\textbf{Procedure:}
\begin{enumerate}
  \item Initialize LIS3DH in Normal mode, $\pm$4g, 100~Hz.
  \item Read WHO\_AM\_I register (0x0F).
  \item \textbf{Expected:} \texttt{0x33} (LIS3DH device ID, datasheet §8.1).
\end{enumerate}

\textbf{Code:}
\begin{lstlisting}[language=C, caption={WHO\_AM\_I verification test case}, 
label={code:testcase1}]
static void testcase1() {
    printf("--- Test Case 1: WHO_AM_I Verification ---\n");
    lis3dh_config.fscale = LIS3DH_FS_4G; 
    lis3dh_config.mode = LIS3DH_MODE_NORMAL;
    // ... init ...
    HAL_StatusTypeDef status = lis3dh_init(&lis3dh_acc, &hi2c1, 
                                          LIS3DH_I2C_DEFAULT_ADDRESS, &lis3dh_config);
    if (status == HAL_OK) 
        printf("PASS: LIS3DH init OK, WHO_AM_I = 0x33\n");
    else 
        printf("FAIL: LIS3DH init failed\n");
}
\end{lstlisting}


%%%%%

\subsubsection{Test Case 2: Operating Mode Configuration}
\textbf{Objective:} Validate that CTRL\_REG1/CTRL\_REG4 correctly configure all 
operating modes.

\textbf{Procedure:} For each mode (Normal/High-Res/Low-Power): configure → read → verify.

\textbf{Code:}
\begin{lstlisting}[language=C, caption={Operating mode configuration test}, 
label={code:testcase2}]
static bool testcase21() {  // Normal mode
    lis3dh_set_mode(&lis3dh_acc, LIS3DH_MODE_NORMAL);
    lis3dh_mode_t mode = lis3dh_get_mode(&lis3dh_acc);
    return (mode == LIS3DH_MODE_NORMAL);
}
// Similar tests for HIGH_RES and LOW_POWER
\end{lstlisting}

%%%%%



\subsubsection{Test Case 3: Full-Scale Configuration}
\textbf{Objective:} Validate configuration of $\pm$2g/$\pm$4g/$\pm$8g/$\pm$16g ranges 
in CTRL\_REG4.

\textbf{Procedure:} For each full-scale range: configure → read → verify.

\textbf{Code:}
\begin{lstlisting}[language=C, caption={Full-scale range configuration test}, 
label={code:testcase3}]
static bool testcase31() {  
    lis3dh_set_full_scale(&lis3dh_acc, LIS3DH_FS_2G);
    lis3dh_fscale_t fs = lis3dh_get_full_scale(&lis3dh_acc);
    return (fs == LIS3DH_FS_2G);
}
// Similar tests for +-4g/+-8g/+-16g
\end{lstlisting}

%%%%%%%%%%%%%%%%%

\subsubsection{Test Case 4: Acceleration Data Injection and Conversion}
\textbf{Objective:} Verify complete data flow: QOM injection → registers → I$^2$C read 
→ firmware scaling.

\textbf{Test sequence:}
\begin{enumerate}
  \item Configure Normal mode, $\pm$4g (So = 8~mg/LSB).
  \item Inject via QEMU: \texttt{qom-set accel-x 1000}.
  \item Firmware reads via \texttt{lis3dh\_read\_acc()} → verifies 1000~mg.
\end{enumerate}

\textbf{Code:}
\begin{lstlisting}[language=C, caption={Acceleration data injection test}, 
label={code:testcase4}]
lis3dh_config.fscale = LIS3DH_FS_4G; 
lis3dh_config.mode = LIS3DH_MODE_NORMAL;
// ... config & enable DRDY ...
// QEMU: qom-set accel-x 1000
lis3dh_read_acc(&lis3dh_acc);
printf("INT! X: %.2f mg", lis3dh_acc.data->acc_mg[0]);  // Expected: 1000.00
\end{lstlisting}

%%%%%%%%%%%%%%%%%%%

\subsubsection{Test Case 5: Data-Ready Interrupt Generation}
\textbf{Objective:} Validate INT1 interrupt generation when new data is available.

\textbf{Configuration:} \texttt{lis3dh\_enable\_drdy\_interrupt()}.

\textbf{Test sequence:}
\begin{enumerate}
  \item Inject X/Y/Z data via QOM properties.
  \item Verify EXTI1 (PA1) triggers → \texttt{HAL\_GPIO\_EXTI\_Callback()}.
  \item Read data + INT1\_SRC (clears latched interrupt).
  \item INT1 returns low.
\end{enumerate}

\textbf{Code:}
\begin{lstlisting}[language=C, caption={Data-Ready interrupt handling}, 
label={code:testcase5}]
void HAL_GPIO_EXTI_Callback(uint16_t GPIO_Pin) {
    if (GPIO_Pin == LIS3DH_INT1_Pin) g_data_ready = true;
}
// Main loop
if (g_data_ready)
{
  g_data_ready = false;  // Clear flag

  // Read acceleration data
  status = lis3dh_read_acc(&lis3dh_acc);
  if (status == HAL_OK)
  {
    printf("INT! X: %.2f mg, Y: %.2f mg, Z: %.2f mg\n",
        lis3dh_acc.data->acc_mg[0],
        lis3dh_acc.data->acc_mg[1],
        lis3dh_acc.data->acc_mg[2]);
  }
  else
  {
    printf("STM32: Read acc ERROR after interrupt\n");
  }

  // Read INT1_SRC to verify/clear interrupt source
  uint8_t int_src = 0;
  lis3dh_read_int_source(&lis3dh_acc, &int_src);
  printf("INT1_SRC: 0x%02X\n", int_src);
}
\end{lstlisting}






























% ==== Test Cases & Results ==== %
\section{Results}

After executing QEMU loaded with the test firmware using the command shown in 
Figure~\ref{code:qemu_test_cmd}:

\begin{figure}[H]
  \begin{lstlisting}[language=bash, caption={QEMU command for automated LIS3DH testing}, label={code:qemu_test_cmd}]
sudo ./qemu-system-arm -M netduinoplus2 -kernel firmware_test_4.elf -serial stdio
  \end{lstlisting}
\end{figure}



the following output confirms that the system fully meets all five testing objectives 
established at the beginning of the chapter:

\begin{lstlisting}[language=bash, caption={Automated test execution output}, 
label={code:test_output}]
=== LIS3DH EMULATION TEST SUITE ===

--- Test Case 1: WHO_AM_I Verification ---
 PASS: LIS3DH init OK, WHO_AM_I = 0x33

--- Test Case 2: Operating Mode Configuration ---

 Test Case 2.1: LIS3DH_MODE_NORMAL
	 TEST PASSED
 Test Case 2.2: LIS3DH_MODE_HIGH_RES
[LIS3DH] Mode: high
[LIS3DH] acc sensitivity: 1.0
	 TEST PASSED
 Test Case 2.3: LIS3DH_MODE_LOW_POWER
[LIS3DH] Mode: low
[LIS3DH] acc sensitivity: 16.0
[LIS3DH] Mode: low
[LIS3DH] acc sensitivity: 16.0
	 TEST PASSED

 PASSED:  Operating Mode Configuration Test case

--- Test Case 3: Full-Scale Configuration ---
 Test Case 3.1: LIS3DH_FS_2G
PASS
 Test Case 3.2: LIS3DH_FS_4G
[LIS3DH] FS: LIS3DH_FS_4G
[LIS3DH] acc sensitivity: 32.0
PASS
 Test Case 3.3: LIS3DH_FS_8G
[LIS3DH] FS: LIS3DH_FS_8G
[LIS3DH] acc sensitivity: 64.0
PASS
 Test Case 3.4: LIS3DH_FS_16G
[LIS3DH] FS: LIS3DH_FS_16G
[LIS3DH] acc sensitivity: 192.0
PASS

 PASS: Full-scale Test case

--- Test Case 4 & 5: Full-Scale Configuration ---
[LIS3DH] Mode: normal
[LIS3DH] acc sensitivity: 48.0
[LIS3DH] FS: LIS3DH_FS_4G
[LIS3DH] acc sensitivity: 8.0
[LIS3DH] ODR: LIS3DH_ODR_100

STM32: LIS3DH initialized with DRDY interrupt on INT1
LIS3DH: Input: 1000.0 mg -> 125.0 LSB -> 125 raw

LIS3DH: out_x_h: 0x1f, out_x_l: 0x40
LIS3DH: Input: 1000.0 mg -> 125.0 LSB -> 125 raw

LIS3DH: out_y_h: 0x1f, out_y_l: 0x40
LIS3DH: Input: 1000.0 mg -> 125.0 LSB -> 125 raw

LIS3DH: out_z_h: 0x1f, out_z_l: 0x40
LIS3DH: INT1 pulsed (DRDY, non-latched)
Calculated So = 8.000 mg/LSB (mode=0, fscale=0)
INT! X: 1000.00 mg, Y: 1000.00 mg, Z: 1000.00 mg
INT1_SRC: 0x00
\end{lstlisting}

\section{Objective-Result Mapping}

\begin{table}[H]
    \centering
    \small
    \caption{Test Objective to Result Mapping}
    \label{tab:test_objective_mapping}
    \begin{tabularx}{\textwidth}{|X|c|c|X|}
        \hline
        \textbf{Objective} & \textbf{Test Case} & \textbf{Result} & \textbf{Justification} \\
        \hline
        Protocol Fidelity & Test Case 1 & \textcolor{green}{PASSED} & 
        I$^2$C subaddressing;  auto-increment correct \\
        \hline
        Configuration Correctness & Test Case 2 \& 3 & \textcolor{green}{PASSED} & 
        Write in the control registers change the internal logic \\
        \hline
        Data Acquisition & Test Case 4 & \textcolor{green}{PASSED} & 
        QOM → registers → perfect mg scaling \\
        \hline
        Interrupt Behavior & Test Case 5 & \textcolor{green}{PASSED} & 
        DRDY → EXTI → clear non-latched working \\
        \hline
        HAL Compatibility & All (Native) & \textcolor{green}{PASSED} & 
        \texttt{HAL\_I2C\_MemRead} \texttt{/Write()} unmodified \\
        \hline
    \end{tabularx}
\end{table}

These tests, though simple in description, together with additional undocumented 
validations, verify the complete functionality of the system developed in this thesis:

\begin{itemize}
  \item QEMU LIS3DH emulator with QOM injection and 
  DRDY interrupt.
  \item STM32F405 I$^2$C controller with datasheet 
  register fidelity and native HAL support.
  \item Interrupt-driven firmware driver, portable 
  and reusable.
  \item End-to-end validation from QOM properties 
  to scaled mg output.
\end{itemize}

The results obtained demonstrate that the emulated system 
behaves identically to its physical counterpart, meeting all established objectives 
and fully validating the technical correctness of the work performed.
